\documentclass[12pt]{article}

\usepackage[margin=1in]{geometry}
\usepackage{setspace}
\usepackage[colorlinks=true,urlcolor=black]{hyperref}

\onehalfspacing

\begin{document}

\section*{Declarations}

\noindent\textbf{Manuscript title:} Government Certification as Technology Governance: Innovation Signals and Institutional Voids in China

\vspace{0.8em}

\noindent\textbf{Author:} Wang Weijia

\vspace{0.8em}

\noindent\textbf{Corresponding author:} Wang Weijia, School of Economics and Finance, Xi'an Jiaotong University, 28 Xianning West Road, Xi'an, Shaanxi 710061, China. Email: \href{mailto:wangweijia@stu.xjtu.edu.cn}{wangweijia@stu.xjtu.edu.cn}.

\vspace{0.8em}

\noindent\textbf{Funding:} The author received no external grant funding for this research.

\vspace{0.8em}

\noindent\textbf{Competing interests:} The author declares no competing financial or personal interests.

\vspace{0.8em}

\noindent\textbf{Ethics approval:} This study uses archival firm-level data and public policy records. It does not involve human subjects, experiments, or personally identifiable survey data.

\vspace{0.8em}

\noindent\textbf{Data availability:} Replication code and processed nonrestricted materials can be made available upon request during review and upon acceptance. Firm-level financial and patent data from CSMAR are subject to the database provider's licensing terms. SRDI certification records are hand-collected from official MIIT and provincial government announcements.

\vspace{0.8em}

\noindent\textbf{Author contribution:} Wang Weijia is the sole author and is responsible for conceptualization, data collection, empirical analysis, writing, and revision.

\vspace{0.8em}

\noindent\textbf{Generative AI disclosure:} AI-assisted language editing was used to improve wording, organization, and journal fit during revision. The author reviewed the manuscript and remains responsible for the accuracy, interpretation, and integrity of all content. In accordance with journal policy, no AI tool is listed as an author.

\section*{Author Biography}

Wang Weijia is affiliated with the School of Economics and Finance at Xi'an Jiaotong University. His research interests include technology governance, innovation policy, financial constraints, and applied econometrics. His current work examines how public certification and related governance instruments shape firm-level technological capability, innovation pathways, and regional development in China.

\end{document}
